\begin{table}
\begin{threeparttable}[htbp]
\centering
\caption{Standardized Effect Sizes}
\label{tab:sde}
\small
\begin{adjustbox}{max width=\textwidth}
\begin{tabular}{lcccccc}
\toprule
Outcome & $\hat{\beta}$ & SE & SD($Y$) & SDE & SE(SDE) & Classification \\
\midrule
Racial Earnings Gap & -0.0284 & 0.0318 & 0.3602 & -0.0788 & 0.0883 & Moderate negative \\
Overall Earnings (White) & 0.0148 & 0.0103 & 0.3602 & 0.0410 & 0.0286 & Small positive \\
Racial Employment Gap & 0.0505 & 0.0469 & 2.2662 & 0.0223 & 0.0207 & Small positive \\
Racial Hiring Gap & 0.1649 & 0.0477 & 1.6593 & 0.0994 & 0.0288 & Moderate positive \\
\bottomrule
\end{tabular}
\end{adjustbox}
\begin{tablenotes}
\small
\item \textit{Notes:} \textbf{Country:} United States. \textbf{Research question:} Did Section 232 steel and aluminum tariffs (25\% and 10\% ad valorem, effective March 2018) narrow or widen the Black-White earnings, employment, and hiring gaps in manufacturing communities with varying tariff exposure? \textbf{Policy mechanism:} Presidential proclamation imposed tariffs on steel (25\%) and aluminum (10\%) imports under national security authority, raising domestic prices of these metals and increasing demand for domestic producers while raising input costs for downstream manufacturing users. \textbf{Outcome definition:} Log average monthly earnings, log beginning-of-quarter employment, and log all hires from the QWI, measuring county-level manufacturing labor market outcomes separately for Black and White workers. \textbf{Treatment:} Continuous --- county-level share of 2016 employment in primary metals (NAICS 331) and fabricated metals (NAICS 332), interacted with post-2018Q2 indicator and Black race indicator. \textbf{Data:} Census QWI Race-Hispanic panel and County Business Patterns 2016, 2015Q1--2020Q1, county-race-quarter level, manufacturing sector (NAICS 31--33). \textbf{Method:} Triple-difference (county $\times$ tariff exposure $\times$ race) with county$\times$race and race$\times$quarter fixed effects; standard errors clustered at state level. \textbf{Sample:} Manufacturing sector in 2,737 counties with positive manufacturing employment in 2016 CBP; 857 high-exposure and 1,880 low-exposure counties; 313,362 county-race-quarter observations for earnings. SDE $= \hat{\beta} / \text{SD}(Y)$ where SD($Y$) is the pre-treatment standard deviation. Classification refers to magnitude, not statistical significance: Large ($|$SDE$| > 0.15$), Moderate ($0.05$--$0.15$), Small ($0.005$--$0.05$), Null ($< 0.005$).
\end{tablenotes}
\end{threeparttable}
\end{table}
